% Hand-edited methodology, 2026-09-05. No empirical results are introduced.
% Operational provenance and notation crosswalk: methodology-notes.md.
\section{Methodology}
\label{sec:methodology}

We compare activation interventions in randomly initialized Pythia-family
models~\citep{biderman2023pythia}, using MiniPile for pretraining and
validation~\citep{kaddour2023minipile}. The methodology separates what an
intervention changes, how its zeros are measured, and how much computation
its placement can affect.

\subsection{Interventions}
\label{sec:interventions}

Each condition specifies a gate-site set $\mathcal{G}$, a gate family and
threshold $\kappa$, and a pressure method with target set $\mathcal{P}$.
Figure~\ref{fig:intervention-ladder} locates the sites: the normalized attention
and FFN inputs ($a,m$), FFN hidden activation ($h$), attention context ($z$),
and query, key, and value operands ($q,k,v$). Query and key gates follow rotary
position encoding (RoPE); a gate at $h$ replaces GELU. Unselected sites retain
their original transformations.

We use ReLU and two fixed-threshold gates:
$G_{+,\kappa}(x)=x\mathbf{1}\{x\geq\kappa\}$ and
$G_{\pm,\kappa}(x)=x\mathbf{1}\{|x|\geq\kappa\}$, with $\kappa\geq0$.
One-sided thresholding removes negative values and small positive values;
symmetric thresholding removes small magnitudes while preserving the signs
and values of survivors. At $\kappa=0$, their forward maps are ReLU and the
identity, respectively. A4 gates $\{a,m,h,z\}$ one-sidedly; A7 adds symmetric
gates at $\{q,k,v\}$. Gate derivatives and exact placements are given in
Appendix~\ref{app:interventions}.

Pressure uses $\mathcal{L}_1$, the equally weighted mean of the mean absolute
activations of each targeted site--layer tensor, measured after gating.
\emph{Naive L1} optimizes $\mathcal{L}_{\mathrm{task}}+\lambda\mathcal{L}_1$.
\emph{Orthogonal L1} (OL1) instead makes a task-only AdamW update, then adds
a pressure correction preconditioned by the task's second-moment estimates.
When this correction opposes the adaptive task direction, OL1 projects out
its conflicting component. A positive budget caps the correction's relative norm.
This controls the update geometry without guaranteeing preserved task loss;
Appendix~\ref{app:pressure} defines the procedure.

Matched contrasts share initialization, data order, and training budget.
Gate and pressure sets are specified independently: A4-OL1 versus A7-OL1
changes both sets, so it compares complete recipes rather than isolating gate
placement at a fixed pressure objective. The ladder organizes these choices;
it does not imply a fully crossed experiment.

\subsection{Quality and model-wide sparsity}
\label{sec:measurement}

We pair validation cross-entropy with sparsity from the same checkpoint and
forward configuration. Local sparsity is the fraction of exactly zero entries
at a named site. To weight zeros by the work they affect, we define
\emph{model-wide sparsity} as $\Smodel=N_0/N_{\mathrm{model}}$, reported as
$100\Smodel\%$. Here, $N_0$ counts scalar multiplications with at least one
zero activation operand in the blocks' QKV, attention-score, probability--value,
attention-output, and two FFN matrix products. $N_{\mathrm{model}}$ counts all
valid multiplications in those operations and the dense language-model head;
the head receives no zero credit. Counts are pooled before division. Thus the
percentage refers to this declared multiplication workload, not all operations
or elapsed time. We count only causally valid pairs in a full, uncached
sequence; Appendix~\ref{app:accounting} gives the exact counters and coverage.

\subsection{Architectural reach}
\label{sec:reach}

The selected sites determine which operations can receive zero operands.
We define their \emph{all-zero reach ceiling},
$\Smax=N_{\mathrm{reach}}(\mathcal{G})/N_{\mathrm{model}}$, by setting those
sites entirely to zero and counting the affected multiplications once.
For example, zeroing $h$ reaches the FFN down projection, whereas zeroing $v$
reaches both the probability--value product and the attention-output projection.
The ceiling depends on architecture and sequence length, but not pressure or
a checkpoint. It explains why identical placements cover different shares of
the workload across model sizes (Figure~\ref{fig:intervention-ladder}).
It is a limit on selected-site reach, not sparsity achievable at acceptable
quality: observed $\Smodel$ also includes natural zeros outside that reach.
The derivation is in Appendix~\ref{app:ceiling}; execution gains require
separate timing.
